
% ==============================================================================
% 1. SISTEMA E CODIFICAÇÃO (Sempre no início)
% ==============================================================================
\usepackage[utf8]{inputenc}       % Suporte para acentuação direta
\usepackage[T1]{fontenc}          % Codificação de fonte correta para português
\usepackage[brazil]{babel}        % Tradução e hifenização em português
\usepackage{indentfirst}          % Indenta o primeiro parágrafo de cada seção
\usepackage{lmodern}              % Fonte Latin Modern (limpa e profissional)
\usepackage{csquotes}             % Gerenciamento avançado de aspas e citações


% ==============================================================================
% 2. MATEMÁTICA E TEOREMAS
% ==============================================================================
\usepackage{amsmath}              % Recursos matemáticos avançados
\usepackage{amsfonts}             % Fontes matemáticas
\usepackage{amssymb}              % Símbolos matemáticos adicionais
\usepackage{amsthm}               % Configuração de ambientes de teoremas

% Configuração dos Teoremas (Traduzidos para o português)
\newtheorem{Axiom}{Axioma}[section]
\newtheorem{Conjecture}{Conjectura}[section]
\newtheorem{Definition}{Definição}[section]
\newtheorem{Theorem}{Teorema}[section]
\newtheorem*{Proof}{Demonstração}
\newtheorem{Example}{Exemplo}[section]
\newtheorem{Corolary}{Corolário}[section]
\newtheorem{Proposition}{Proposição}[section]

% ==============================================================================
% 3. GRÁFICOS, CORES E DIAGRAMAÇÃO
% ==============================================================================
\usepackage[top=3cm, bottom=2cm, left=3cm, right=2cm]{geometry} % Margens padrão
\usepackage{graphicx}             % Suporte para inserção de imagens e figuras
\graphicspath{{imgs/}}         % Define a pasta padrão para buscar fotos/gráficos
\usepackage{microtype}            % Melhora o espaçamento entre letras e evita hífens excessivos
\usepackage{fancyhdr}             % Permite customizar cabeçalhos e rodapés
\usepackage{float}                % Controle estrito de posição de elementos ([H])
\usepackage{pgfplots}             % Criação de gráficos diretamente no LaTeX
\pgfplotsset{compat=1.18}         % Boa prática: define a compatibilidade do pgfplots

% ==============================================================================
% 4. TABELAS AVANÇADAS
% ==============================================================================
\usepackage{booktabs}             % Linhas de tabela profissionais (\toprule, etc)
\usepackage{multirow}             % Linhas múltiplas em tabelas
\usepackage{tabularx}             % Tabelas com largura ajustável automática
\usepackage{makecell}             % Quebras de linha e alinhamento dentro de células

% ==============================================================================
% 5. CÓDIGOS DE PROGRAMAÇÃO
% ==============================================================================
\usepackage{listings}             % Exibição de blocos de código fonte

% ==============================================================================
% 6. LINKS, URLs E REFERÊNCIAS (Sempre por último)
% ==============================================================================
\usepackage{url}                  % Formatação correta de links de internet
\usepackage[num]{abntex2cite}     % Normas ABNT para citações
\usepackage{hyperref}             % Transforma sumário, links e citações em hiperlinks clicáveis
\hypersetup{
	colorlinks=true,              % Ativa cores nos links em vez de caixas coloridas
	linkcolor=blue,               % Cor dos links internos (sumário, figuras)
	citecolor=green,              % Cor das citações bibliográficas
	filecolor=magenta,            
	urlcolor=cyan                 % Cor de links de sites externos
}

%\usepackage{glossaries}
